\documentclass[12pt,a4paper]{article}

% Encoding and fonts
\usepackage[utf8]{inputenc}
\usepackage[T1]{fontenc}
\usepackage{lmodern}
\usepackage{textcomp}

% Page layout
\usepackage[top=1in, bottom=1in, left=1in, right=1in]{geometry}

% Typography
\usepackage{microtype}
\usepackage{parskip}

% Language
\usepackage[english]{babel}

% Headers and footers
\usepackage{fancyhdr}
\pagestyle{fancy}
\fancyhf{}
\fancyhead[L]{\leftmark}
\fancyhead[R]{\thepage}
\fancyfoot[C]{\thepage}
\renewcommand{\headrulewidth}{0.4pt}
\renewcommand{\footrulewidth}{0pt}

% Section styling
\usepackage{titlesec}
\titleformat{\section}{\Large\bfseries}{\thesection}{1em}{}
\titleformat{\subsection}{\large\bfseries}{\thesubsection}{1em}{}
\titleformat{\subsubsection}{\normalsize\bfseries}{\thesubsubsection}{1em}{}
\titlespacing*{\section}{0pt}{2.5ex plus 1ex minus .2ex}{1.5ex plus .2ex}
\titlespacing*{\subsection}{0pt}{2ex plus 1ex minus .2ex}{1ex plus .2ex}
\titlespacing*{\subsubsection}{0pt}{1.5ex plus 1ex minus .2ex}{0.8ex plus .2ex}

% List formatting
\usepackage{enumitem}
\setlist[itemize]{topsep=0.5ex, itemsep=0.25ex, parsep=0pt}
\setlist[enumerate]{topsep=0.5ex, itemsep=0.25ex, parsep=0pt}

% Essential packages
\usepackage{graphicx}
\usepackage[table]{xcolor}
\usepackage{booktabs}
\usepackage{array}
\usepackage{tabularx}
\usepackage{longtable}
\usepackage{amsmath}
\usepackage{amssymb}
\usepackage[normalem]{ulem}
\rowcolors{2}{gray!10}{white}
\renewcommand{\arraystretch}{1.3}

% Cross-references
\usepackage{hyperref}
\usepackage{cleveref}

% Custom commands
\newcommand{\pilcrow}{\S}

% Hyperref setup
\hypersetup{
    colorlinks=true,
    linkcolor=blue,
    filecolor=magenta,
    urlcolor=cyan,
}

% Code listings
\usepackage{listings}
\definecolor{codebg}{RGB}{248,248,248}
\definecolor{codeframe}{RGB}{200,200,200}
\definecolor{codekeyword}{RGB}{0,0,180}
\definecolor{codecomment}{RGB}{0,128,0}
\definecolor{codestring}{RGB}{163,21,21}
\lstset{
    basicstyle=\ttfamily\small,
    breaklines=true,
    breakatwhitespace=false,
    frame=single,
    framerule=0.5pt,
    rulecolor=\color{codeframe},
    backgroundcolor=\color{codebg},
    keepspaces=true,
    numbers=left,
    numberstyle=\tiny\color{gray},
    numbersep=8pt,
    showstringspaces=false,
    tabsize=4,
    xleftmargin=1em,
    xrightmargin=1em,
    aboveskip=1em,
    belowskip=1em,
    keywordstyle=\color{codekeyword}\bfseries,
    commentstyle=\color{codecomment}\itshape,
    stringstyle=\color{codestring},
    literate={_}{{\textunderscore}}1
             {-}{{\textminus}}1
             {<}{{\textless}}1
             {>}{{\textgreater}}1
             {~}{{\textasciitilde}}1
             {^}{{\textasciicircum}}1
}

% YAML language definition for listings
\lstdefinelanguage{YAML}{
    keywords={true,false,null,y,n},
    keywordstyle=\color{blue}\bfseries,
    sensitive=false,
    comment=[l]{\#},
    morecomment=[s]{/*}{*/},
    commentstyle=\color{gray}\ttfamily,
    stringstyle=\color{red}\ttfamily,
    morestring=[b]'',
    morestring=[b]'',
}

% TOML language definition for listings
\lstdefinelanguage{TOML}{
    keywords={true,false},
    keywordstyle=\color{blue}\bfseries,
    sensitive=true,
    comment=[l]{\#},
    commentstyle=\color{gray}\ttfamily,
    stringstyle=\color{red}\ttfamily,
    morestring=[b]'',
    morestring=[b]'',
}

% Dockerfile language definition for listings
\lstdefinelanguage{Dockerfile}{
    keywords={FROM,RUN,CMD,LABEL,EXPOSE,ENV,ADD,COPY,ENTRYPOINT,VOLUME,USER,WORKDIR,ARG,ONBUILD,STOPSIGNAL,HEALTHCHECK,SHELL},
    keywordstyle=\color{blue}\bfseries,
    sensitive=false,
    comment=[l]{\#},
    commentstyle=\color{gray}\ttfamily,
    stringstyle=\color{red}\ttfamily,
    morestring=[b]'',
    morestring=[b]'',
}

% JSON language definition for listings
\lstdefinelanguage{JSON}{
    keywords={true,false,null},
    keywordstyle=\color{blue}\bfseries,
    sensitive=true,
    commentstyle=\color{gray}\ttfamily,
    stringstyle=\color{red}\ttfamily,
    morestring=[b]'',
}

% Code listing float environment
\usepackage{float}
\newfloat{listing}{htbp}{lol}[section]
\floatname{listing}{Listing}

% Admonition boxes
\usepackage{tcolorbox}
\tcbuselibrary{skins,breakable}

% Admonition style definitions
\newtcolorbox{admonitionbox}[2][]{
    enhanced,
    breakable,
    colback=#2!5,
    colframe=#2!80!black,
    fonttitle=\bfseries,
    title=#1,
    left=2mm,
    right=2mm,
}

% Synthon tuple boxes
\newtcolorbox{synthonbox}{
    enhanced,
    colback=blue!5,
    colframe=blue!75!black,
    fontupper=\large,
    center upper,
    boxrule=1pt,
    arc=4pt,
}

\begin{document}

\title{Structural Probing of the EML Sheffer Operator via the SynthOmnicon Grammar}
\date{2026-04-21}

\maketitle

\textbf{Keywords:} EML, Sheffer operator, elementary functions, SynthOmnicon, Frobenius, symbolic regression

\subsection{Abstract}

There is a fact about Odrzywołek''s EML Sheffer operator that his paper cannot
state, because it requires a vocabulary the paper does not have. The operator
$\text{eml}(x,y) = e^x - \ln y$, paired with the constant $1$, is not merely a
compact basis for elementary functions --- it is the algebraically forced ceiling
of a structural tier, the highest sub-Frobenius type in the elementary function
algebra. Every feature of the discovery --- why subtraction and not addition, why
the terminal must be $1$, why three variants exist and cannot be collapsed into
one, why gradient recovery succeeds at depth 4 and fails at depth 6, why complex
arithmetic is mandatory --- follows from a single structural fact: EML encodes the
exp/ln duality as an assertion rather than a proof. Closing that gap requires not
a better search but a different algebraic kind.

We apply the grammar\textbackslash{}footnote\{SynthOmnicon 12-primitive grammar. Source: \textbackslash{}url\{https://github.com/umpolungfish/synthomnicon\_private\}\} to make this precise. It is not being used here as an
oracle; it is being used as a coordinate system that
forces questions the paper does not ask. Several of the answers are surprising
even from within the grammar. One of them --- that the paper''s own ternary candidate
$T(x,y,z)$ is Calc 1 dressed in ternary notation --- is a case where the grammar
reads the paper against itself. The paper, to its credit, supplies its own
refutation two paragraphs before the acknowledgments.

\begin{center}
\rule{0.5\textwidth}{0.4pt}
\end{center}

\subsection{I. The Structural Type of EML}

Any encoding carries the risk of projecting a framework onto material that does
not fit, so the first obligation is to earn the type assignment. The EML operator
encodes:

\[
\mathbf{eml} = \langle D_{\infty};\ T_{\bowtie};\ R_\text{dagger};\ P_{\pm};\
F_{\hbar};\ K_\text{slow};\ G_{\aleph};\ \Gamma_\text{seq};\ \Phi_c;\ H_1;\
S_{1:1};\ \Omega_Z \rangle
\]

Most coordinates are uncontroversial. The assignments that require argument are
$\Phi_c$, $T_{\bowtie}$, and $P_{\pm}$, in that order of difficulty.

$\Phi_c$ in a functional algebra rather than a physical system is not an obvious
assignment. The critical manifold here is the snapping manifold identified in \pilcrow{}VI:
the set of parameter vectors in the master formula that correspond exactly to
symbolic EML trees. The snapping phenomenon --- gradient weights collapsing to
binary values --- is a criticality transition in the parameter space, and its
existence at all depths (perturbed-truth initialization always recovers) is the
non-trivial fact that earns $\Phi_c$.

$T_{\bowtie}$ is assigned not to the topology of the elementary functions themselves
but to the structural relationship between the grammar analysis and the paper.
The two threads --- Odrzywołek''s empirical scaffolding and the grammar''s algebraic
diagnosis --- cross and exchange at \pilcrow{}VII, where the paper''s own comparison to
quantum computing turns out to be structurally exact in a sense the author did
not intend.

$P_{\pm}$ is argued in \pilcrow{}II and \pilcrow{}IV. The short version: EML''s exp/ln self-duality is
real but external --- to access the symmetric cousin $-\text{eml}(y,x)$ you must
leave the algebra. $P_{\pm}^{\text{sym}}$ requires the duality to be internal,
and it is not.

This type is $O_2^\dagger$ by tier rule R5: $\Phi_c$ present, $\Omega_Z \neq
\Omega_0$, and $D_{\infty}$.

\begin{center}
\rule{0.5\textwidth}{0.4pt}
\end{center}

\subsection{II. Why Subtraction and Not Addition or Multiplication}

The wrong answer first, because it is instructive --- and because it is also the
paper''s answer, offered without knowing it. The obvious candidate for why
subtraction succeeds is asymmetry: subtraction is non-commutative, and
non-commutativity gives the tree a chirality, which gives the grammar inversion
capabilities. But this is not the reason. The operator $B(x,y) = x - y/2$ is
equally non-commutative and fails --- the paper''s own \pilcrow{}5 gives this counterexample
explicitly, noting that $B(B(x,x), x) = 0$ but $B$ still cannot generate all
elementary functions. Asymmetry is necessary but not sufficient.

The grammar''s answer is about the relational mode primitive $R$. The two inputs of
$\text{eml}(x,y)$ do not meet as equals: $x$ has already been transformed by
$\exp$ (the homomorphism $(\mathbb{R},+) \to (\mathbb{R}^+,\times)$) and $y$ has
been transformed by $\ln$ (the reverse). Subtraction then computes the gap between
an element of the multiplicative group and a multiplicative-group element
re-expressed in additive coordinates. This is $R_\text{dagger}$ --- the adjoint
relational mode. The operation measures how far $x$''s image under $\exp$ departs
from $y$''s pre-image under $\ln$; it is directed, oriented, and asymmetric
precisely because it is the adjoint of the bridge between two groups, not merely
a difference of two arbitrary quantities.

Addition would collapse this into $R_\text{super}$ (symmetric inclusion of both
inputs), losing the directed structure. Multiplication would give $R_\text{cat}$
(each factor participates catalytically, staying within the multiplicative group),
losing the additive-to-multiplicative bridge. Division, as EDL demonstrates,
retains $R_\text{dagger}$ but encodes a ratio in the multiplicative group rather
than a translated difference --- a valid variant but not the primary one, for reasons
\pilcrow{}III will make clear.

The parity assignment $P_{\pm}$ follows. EML knows the distinction between its
arguments --- $\text{eml}(x,y) \neq \text{eml}(y,x)$ --- but the symmetry between
them is accessible only externally: you must negate and swap to reach
$-\text{eml}(y,x)$. For $P_{\pm}^{\text{sym}}$, the Frobenius parity, the
symmetry must be internal: $\mu \circ \delta = \text{id}$ within the same algebra.
This distinction --- assertion of duality versus proof of duality --- is the structural
fact the paper cannot name and cannot reach by any extension of its own methods.

\begin{center}
\rule{0.5\textwidth}{0.4pt}
\end{center}

\subsection{III. Why the Terminal Constant is Forced to be $1$}

This follows from \pilcrow{}II, though it is not obvious that it does.

If the adjoint structure of $R_\text{dagger}$ is what makes EML work, then the
terminal symbol must be the element that makes the adjoint degenerate gracefully
--- the point where one arm of the bridge can be switched off without destroying
the other. That point is $1$. It is the unique element where the additive and
multiplicative group identities coincide: $\ln(1) = 0$ (additive neutral) and
$1 = e^0$ (return to multiplicative neutral). Plugging $y = 1$:


\[\text{eml}(x, 1) = e^x - \ln 1 = e^x - 0 = e^x,\]


the $\ln$ arm is annihilated by the additive identity, and the operator reduces
to pure $\exp$. No other constant has this property. $\ln(e) = 1 \neq 0$; adding
$1$ to $e^x$ gives $e^x + 1$, not $e^x$.

The grammar reads $1$ as the $\Omega_0$ sector of the algebra: the terminal with
trivial winding, no topological cost, the basepoint. Plugging into the $\exp$ slot
instead: $\text{eml}(1, y) = e - \ln y$, with $\text{eml}(1,1) = e$. So $e$ is
not an independent constant but the one-step image of $1$ under EML. The triple
of terminals $\{1, e, -\infty\}$ for the three variants is itself an EML orbit:
$1$ is the monad unit, $e$ is its first iterate, and $-\infty = \ln(0)$ is the
boundary degeneration as $y \to 0^+$. Of the three, $-\infty$ is the most
forthright: it makes no pretense of being a number.

This is the sense in which the three variants are projections of the same
algebraic object, and it answers a question the paper does not ask: why does EDL
require $e$ specifically rather than some other constant? Because $e$ is the image
of $1$ under the operator, and EDL''s division structure places the multiplicative
identity of the output in the role that EML''s $1$ plays additively.

\begin{center}
\rule{0.5\textwidth}{0.4pt}
\end{center}

\subsection{IV. The EML Family as a $\mathbb{Z}_2$ Orbit Stuck Below $O_{\infty}$}

The paper discovers the three variants on different timescales --- EML first, the
cousins ''''a month later'''' --- and treats them as independent discoveries noting they
are ''''close cousins.'''' The grammar sees the orbit structure immediately. Mathematics
is not always patient with timescales, but the interesting question is not whether
the orbit can be named but what it implies.

The three variants --- EML (constant $1$), EDL = $e^x/\ln y$ (constant $e$), and
$-\text{eml}(y,x) = \ln x - e^y$ (constant $-\infty$) --- form an orbit under the
external $\mathbb{Z}_2$ action: swap arguments and negate. The word ''''external'''' is
doing all the work. The orbit is real, the symmetry is real, but it is not
accessible from within the algebra. To go from EML to $-\text{eml}(y,x)$ you
cannot write a finite EML tree --- you must leave the grammar $S \to 1 \mid
\text{eml}(S,S)$ entirely, negate, and relabel. This is less a mathematical
operation than a change of address.

This is the exact content of $P_{\pm}$ rather than $P_{\pm}^{\text{sym}}$: the
duality exists but is external. And it is why $O_{\infty}$ is unreachable from
EML. The Frobenius tier rule R1 requires $\Phi_c$ and $P_{\pm}^{\text{sym}}$
together; EML has $\Phi_c$ but sits at $P_{\pm}$. By Frobenius non-synthesizability
(\pilcrow{}23), no composition of sub-Frobenius factors can produce $P_{\pm}^{\text{sym}}$,
and $P$ is a $\min$-primitive under $\otimes$: stacking EML trees contributes
$P_{\pm}$ at every step, and the bottleneck never lifts.

The implication for Table 2 of the paper is sharp. The ''''?'''' binary at count 2 ---
the conjectured self-bootstrapping operator requiring no distinguished constant ---
must live at $P_{\pm}^{\text{sym}}$. It cannot be approached along the EML
reduction ladder. The author, not having this vocabulary, frames the open question
as a search problem: ''''whether an EML-type binary Sheffer working without pairing
with a distinguished constant exists.'''' The grammar frames it differently: the
question is whether the elementary function algebra admits a Frobenius generator.
Those are the same question, but the second framing shows why the answer cannot
come from better search alone.

\begin{center}
\rule{0.5\textwidth}{0.4pt}
\end{center}

\subsection{V. The Reduction Ladder as a Lattice Meet in $\Gamma$}

The ablation sequence $36 \to 7 \to 6 \to 4 \to 4 \to 3 \to 3$ is presented as
a historical reduction and is exactly that, but it has a simultaneous algebraic
reading: it is a sequence of meets in the interaction grammar primitive $\Gamma$.
The repeated $4$ in the sequence --- Calc 1 and Calc 2 each have four primitives ---
is not a transcription error. They have the same cardinality and completely
different algebraic structure, which is the kind of thing that makes combinatorics
an unreliable guide to understanding.

Calc 3 operates in $\Gamma_\text{and}$: six primitives run in parallel, each
independently callable. The passage to Calc 2 --- dropping negation and reciprocal,
keeping $\exp$, $\ln$, and subtraction --- is a meet that eliminates the
parallelism: negation and reciprocal are now recoverable only through sequential
composition (e.g. $x - x = 0$, $e^0 = 1$, $0 - 1 = -1$). The topology changes
from $\Gamma_\text{and}$ to $\Gamma_\text{seq}$. The passage from Calc 2 to EML
is a further $\Gamma_\text{seq}$ compression: the sequential structure already
present in the expressions $e^x - \ln y$ is absorbed into the operator''s
definition, so that a single call to $\text{eml}$ performs what was previously a
pipeline of $\exp$, $\ln$, and $-$.

What the paper calls ''''generating constants on its own'''' --- the property Calc 3 and
Calc 2 have but Calc 1 lacks --- corresponds to $G = G_{\aleph}$ in the grammar
(global scope: the system bootstraps arbitrary constants from interaction alone).
Calc 1 falls to $G = G_{\gimel}$ because $x^y$ and $\log_x y$ are purely
multiplicative-group operations that cannot generate additive-group elements from
their own iteration without a starting constant. EML recovers $G_{\aleph}$ via the
$\ln(1) = 0$ fixed point, which is why the terminal is forced to be $1$ and not a
general real number.

The ''''?'''' binary at count 2 would require $G_{\aleph}$ with $\Phi_c$ and
$P_{\pm}^{\text{sym}}$ together --- a self-bootstrapping Frobenius generator. The
grammar has no guarantee such an object exists in the class of elementary
functions, and \pilcrow{}IX gives reason to think it does not.

\begin{center}
\rule{0.5\textwidth}{0.4pt}
\end{center}

\subsection{VI. The Snapping Phenomenon is a $\Phi_c$ Transition}

The gradient recovery experiments in \pilcrow{}4.3 of the paper yield a characteristic
depth-dependent success rate: 100\% at depth 2, approximately 25\% at depths 3--4,
below 1\% at depth 5, and zero in 448 attempts at depth 6 (the specificity of
448 suggests either a principled stopping criterion or a very long weekend).
The paper attributes this to ''''basins of attraction'''' becoming hard to find. This is correct but
describes the symptom rather than the cause.

The master formula at depth $n$ has $5 \times 2^n - 6$ continuous parameters.
The snapping target is a discrete set of exact-symbolic parameter vectors, each
a point on a measure-zero submanifold. The $\Phi_c$ manifold --- the submanifold
of symbolic trees --- is present at every depth: this is confirmed directly by the
paper''s own observation that ''''when the weights of the correct EML tree are
perturbed by Gaussian noise, the optimization converges back to the exact values
in 100\% of runs, even for trees of depth 5 and 6.'''' The manifold does not shrink
or disappear. What changes is the ratio of the basin''s measure to the total
parameter volume, which decreases as $2^{-n}$ in the relevant directions.

The harder point --- and this is where the analogy to physical criticality has real
traction --- is why the basin measure decreases rather than stabilizes. The EML
master formula at depth $n$ has $2^n$ discrete symbolic targets but $5 \times
2^n$ continuous parameters per target, and the basin of attraction around each
target is bounded by the separation from neighboring targets in the direction of
steepest descent. As $n$ grows, neighboring targets come closer together in
parameter space (the tree space densifies), and the basins narrow. At depth $\geq 5$,
gradient descent initialized randomly lands in a shallow non-critical basin with
probability approaching 1, and the kinetics --- $K_\text{fast}$ for Adam with
standard step sizes --- cannot escape. This is the $K_\text{trap}$ failure mode:
not that the $\Phi_c$ manifold is absent but that the dynamics cannot actualize
it from a generic starting point.

The grammar''s prediction for the open question in \pilcrow{}4.3 (''''possibly using another
binary operator similar to (4) but with better properties'''') is specific: the
required property is a $\Phi_c$ manifold with larger basin-to-volume ratio at
depth $n$. This is a structural property of the operator''s type, not an
optimization property of the training procedure. An operator with
$P_{\pm}^{\text{sym}}$ would have its snapping targets separated by algebraic
relations that enforce larger basins --- the Frobenius condition creates additional
gradient-landscape symmetry that EML''s $P_{\pm}$ lacks.

\begin{center}
\rule{0.5\textwidth}{0.4pt}
\end{center}

\subsection{VII. Complex Arithmetic as Imscriptive Structure}

There is a sentence in \pilcrow{}5 of the paper that the grammar reads as more exact than
the author intended. Discussing the mandatory complex intermediates, Odrzywołek
writes: ''''Just as quantum computing uses complex amplitudes to compute real
probabilities, EML uses complex intermediates to compute real elementary
functions.'''' He offers this as an analogy, a rhetorical softening of what he calls
a ''''minor problem,'''' and immediately moves on. It is, in fact, the most
structurally significant observation in the paper. The author has discovered
something true and apologized for it. The grammar finds it is not an analogy.
It is a structural identity.

Quantum computing''s complex amplitudes carry winding information --- the phases
$e^{i\theta}$ encode which paths have been traversed --- and the real probabilities
emerge by squaring moduli, which projects the winding information out while
preserving its effects on interference. EML''s complex intermediates carry exactly
the same structure: the imaginary part of $\ln(x)$ for $x < 0$ is $i\pi$,
encoding one unit of winding around the branch point at $0$, and the real
trigonometric output emerges when this imaginary part contributes to the final
$\exp$ through Euler''s formula $e^{i\theta} = \cos\theta + i\sin\theta$. The
mechanism is formally the same: complex bulk encodes winding, real boundary reads
out a projection.

The grammar reads this as $D_{\odot}$ (imscriptive): the real elementary functions
on the boundary are encoded in the complex bulk computation. This is not a
metaphor. The Riemann surface of $\ln$ is literally a bulk: the principal branch
is one sheet, the universal cover contains all of them, and the real axis is a
boundary on each sheet. EML computes on one sheet ($H_1$ chirality --- the choice
of principal branch) and produces real output by projection.

Why is this inevitable? Because $\pi_1(\mathbb{C}^\times) = \mathbb{Z}$. The
fundamental group of the domain of $\ln$ is non-trivial, and any function
inheriting its structure carries $\Omega_Z$ winding. Trigonometric functions have
$\Omega_Z$ in their monodromy, and any operator computing them must pass through
the bulk where this winding lives. A real-only Sheffer --- an operator that computes
trigonometric functions without ever visiting the complex plane --- would require
$\Omega_0$ winding, which is structurally incompatible with the monodromy of
$\sin$ and $\cos$. The author''s search for ''''alternatives using pairs of
trigonometric/hyperbolic functions'''' found nothing because those functions carry
the same $\Omega_Z$ and merely repackage the same bulk structure.

The branch-cut correction the paper handles manually --- the $2\pi i$ jump on the
negative real axis causing the wrong sign for $i$ --- is the $H_1$ chirality
becoming visible as a choice. Two valid chirality orientations exist (principal
branch and its mirror), and EML forces the author to pick one explicitly. Working
on the Riemann surface of $\ln$ proper would be $H_{\infty}$; EML operates at $H_1$
by fixing the chirality and correcting deviations at the cut.

\begin{center}
\rule{0.5\textwidth}{0.4pt}
\end{center}

\subsection{VIII. The EML Closure Boundary}

The paper does not ask which mathematical objects lie outside the EML closure.
The grammar gives a sharp answer for most cases and an honest uncertainty for one.

The sharp cases: everything in the EML closure sits at or below $P_{\pm}$, because
EML contributes $P_{\pm}$ at every composition step and $P$ is a $\min$-primitive
under $\otimes$. By Frobenius non-synthesizability (\pilcrow{}23), $P_{\pm}^{\text{sym}}$
cannot be built from sub-Frobenius factors. The following objects require
$P_{\pm}^{\text{sym}}$:

\textbf{Stark units.} The Shimura reciprocity law equates the Galois action on torsion
points with the CM automorphism. This equation is the Frobenius condition:
$\mu \circ \delta = \text{id}$ at the level of the ring class field algebra. Stark
units, as explicit CM values, carry this condition in their algebraic integrality.
They cannot be written as values of any EML tree.

\textbf{SIC-POVM fiducials.} The minimal polynomial of the fiducial vector is
self-reciprocal --- its coefficient sequence reads the same forwards and backwards.
This palindromic property --- the coefficient sequence reads the same in either
direction, like a good theorem --- is precisely $P_{\pm}^{\text{sym}}$ at the polynomial
level, and it forces the fiducial coordinates into the $(-1)$-eigenspace of the
relevant involution on the unit group. The argument is developed more fully in
SIC\_OQ\_MANUSCRIPT \pilcrow{}IV.

\textbf{The modular $j$-function.} The two functional equations $j(\tau+1) = j(\tau)$
and $j(-1/\tau) = j(\tau)$ are the multiplication and comultiplication legs of
the Frobenius condition on the modular group algebra. They encode together
$\mu \circ \delta = \text{id}$, which is why $j$ is not an elementary function.

The case requiring honest uncertainty is the \textbf{Riemann $\xi$ function}. If the
Riemann hypothesis holds, the zeros lie on $\text{Re}(s) = 1/2$, which would
make the functional equation $\xi(s) = \xi(1-s)$ an exact self-duality at the
$\Phi_c$ manifold --- the grammar''s assignment would be $P_{\pm}^{\text{sym}}$, and
$\xi$ would be outside the EML closure. But the grammar cannot assert the Riemann
hypothesis. What we can say is: if $\xi \in P_{\pm}^{\text{sym}}$, then the
Frobenius condition and the Riemann hypothesis are the same statement at different
coordinate scales. Whether they are, we do not know. We are in good company.

\begin{center}
\rule{0.5\textwidth}{0.4pt}
\end{center}

\subsection{IX. The Ternary Candidate is Structurally Misidentified}

The paper''s proposed ternary:


\[T(x,y,z) = \frac{e^x}{\ln x} \times \frac{\ln z}{e^y}.\]


Simplify: $T(x,y,z) = e^{x-y} \cdot \log_x z$. Setting $y = x$:
$T(x,x,z) = \log_x z$. Setting $z = e^y$: $T(x,y,e^y) = e^{x-y} \cdot y/\ln x$.
This is the Calc 1 system --- $x^y$ and $\log_x y$ --- packed into ternary notation.
The algebraic content is not new; the ternary presentation is. Odrzywołek has
put Calc 1 in a trenchcoat.

This is not to dismiss the ternary route. The paper correctly identifies that a
ternary operator with $T(x,x,x) = c$ for some constant $c$ would have the
potential to generate its own constants from iteration, bypassing the need for a
distinguished terminal. And the grammar agrees: a genuinely Frobenius ternary
operator would satisfy something like $T(T(x,y,z), y, z) = x$ (the Frobenius
condition written as self-invertibility), which would be a structural advance
over EML. But $T$ as presented is not on that path, because it is Calc 1 and
Calc 1 is $G = G_{\gimel}$ --- it requires a distinguished constant ($e$ or $\pi$)
and cannot bootstrap from arbitrary input.

The paper''s own counterexample $B(x,y) = x - y/2$ --- given to illustrate why the
search for a self-bootstrapping binary operator is non-trivial --- applies here too.
$T(x,x,x) = 1$ is the analogue of $B(B(x,x),x) = 0$: a useful algebraic property
that looks like self-bootstrapping but is not. The genuine test is whether $T$ can
recover arbitrary initial conditions from its own iteration, and the Calc 1 algebra
cannot.

\begin{center}
\rule{0.5\textwidth}{0.4pt}
\end{center}

\subsection{X. The Ceiling and the Open Question}

The structural type of EML is:

\[
\mathbf{eml} = \langle D_{\infty};\ T_{\bowtie};\ R_\text{dagger};\ P_{\pm};\
F_{\hbar};\ K_\text{slow};\ G_{\aleph};\ \Gamma_\text{seq};\ \Phi_c;\ H_1;\
S_{1:1};\ \Omega_Z \rangle
\]

Every coordinate arrived under constraint: $R_\text{dagger}$ because the adjoint
structure of subtraction is the only operation that bridges the two group
homomorphisms without collapsing their asymmetry; $P_{\pm}$ because the exp/ln
self-duality is accessible only externally; $\Phi_c$ because the snapping manifold
is present and actualized at accessible depths; $H_1$ because the principal branch
choice is a forced chirality orientation; $\Omega_Z$ because the fundamental group
of the logarithm''s domain is non-trivial. The type is not assigned --- it is derived.

What the type says is that EML is the maximum of a tier, not a point on an
ascending sequence. The reduction ladder cannot continue past EML within the
same algebraic kind. The ''''?'''' binary at count 2 in Table 2 --- the conjectured
self-bootstrapping operator --- cannot be found by pushing the ablation further.
It requires $P_{\pm}^{\text{sym}}$, which is structurally separated from $P_{\pm}$
by Frobenius non-synthesizability. The ladder has reached its ceiling; the next
floor is a different building.

Whether that next floor exists --- whether there is a binary or ternary elementary
function that achieves Frobenius self-duality internally --- is genuinely open.
The grammar cannot resolve it, only sharpen it. The question the paper frames
as ''''does an EML-type Sheffer without a distinguished constant exist?'''' is, under
the grammar''s coordinates, the question: ''''does the class of elementary functions
contain a Frobenius generator?'''' That is a question about the internal symmetry
group of an algebra that mathematicians have been circling for decades from
other directions --- the same question that Stark units, SIC-POVMs, and the
zeros of $\xi$ are each partial answers to, in their respective domains.

If the answer is yes, EML will have been the sign pointing toward it. If no,
EML is already the deepest thing in its tier, and that is enough.

\begin{center}
\rule{0.5\textwidth}{0.4pt}
\end{center}

\subsection{XI. Validation of the Encoding}

\pilcrow{}I opened with an obligation: ''''any encoding carries the risk of projecting a
framework onto material that does not fit, so the first obligation is to earn the
type assignment.'''' What follows is the attempt to discharge that obligation.
ENCODING\_EPISTEMOLOGY \pilcrow{}Summary specifies nine convergence criteria. They are not
a checklist; they are a set of stresses applied from different directions, and
an encoding that survives all of them simultaneously is a finding rather than a
choice.

\textbf{Tier consistency.} $O_2^\dagger$ predicts: a self-referential loop that is
topologically protected and recursively unbounded, but \emph{not} Frobenius-closed ---
encoding and decoding are not mutually inverse. The grammar $S \to 1 \mid
\text{eml}(S,S)$ supplies the loop; the $\Omega_Z$ branch structure supplies
topological protection; $D_{\infty}$ supplies unbounded depth. The non-$O_{\infty}$
prediction is the important one: it requires that no EML tree can recover both
inputs from the output, and the paper never defines an inverse operator because
none exists. The tier is consistent with every known property of the system and
falsified by none. $\checkmark$

\textbf{Nearest-neighbor coherence.} The most structurally unexpected output of an
encoding is its neighborhood, because the encoder cannot control who shows up.
On the shared coordinates $T_{\bowtie}$, $P_{\pm}$, $\Phi_c$, $\Omega_Z$,
$\Gamma_\text{seq}$, the closest catalog neighbor is \emph{The Thunder: Perfect Mind}
--- a text that generates all its content from a single recursive operator asserting
rather than proving its self-duality. That a Nag Hammadi tractate and a symbolic
computation primitive are structural neighbors was not anticipated when the
encoding was made. Also in the neighborhood: conformal field theories at
criticality and the Riemann $\zeta$ function (same coordinates, $P_{\pm}$ pending
the Riemann hypothesis). These are not implausible neighbors. They are
illuminating ones. $\checkmark$

\textbf{Tensor sanity.} This check carries independent weight because it tests a
claim made in \pilcrow{}VIII --- that EML cannot compute SIC-POVM fiducials --- by a
completely different route. $\text{eml} \otimes \text{SIC-POVM}$: the $P$
bottleneck gives $\min(P_{\pm}, P_{\pm}^{\text{sym}}) = P_{\pm}$, pulling the
composite down from $O_{\infty}$. EML is not deficient in $\Phi$ or $\Omega$;
it is deficient in $P$, and the bottleneck rule is ruthless: no additional EML
structure compensates. That \pilcrow{}VIII''s argument and the tensor calculation arrive
at the same structural fact by independent means is the kind of convergence that
matters. $\text{eml} \otimes \text{eml}$ closes on EML''s own type --- the grammar
is closed under self-composition as it must be. $\checkmark$

\textbf{Meet/join coherence.} $\text{eml} \wedge \text{SIC-POVM}$ returns EML''s own
type: EML is the exact common subalgebra of itself and the SIC-POVM, which is
EML plus the Frobenius condition plus one additional chirality level. The join
returns the SIC-POVM type --- the minimal envelope. EML is the floor, the SIC-POVM
is the ceiling, and the gap between them is precisely the Frobenius condition.
The lattice picture is clean and requires no special pleading. $\checkmark$

\textbf{Le Chatelier inversion.} If the equilibrium substrate is not immediately
recognizable, the encoding is suspect. Here it is recognizable: the polynomial
algebra over $\mathbb{R}$,


\[X^* \approx \langle D_{\triangle};\ T_\text{network};\ R_\text{super};\
P_\text{asym};\ F_{\ell};\ K_\text{slow};\ G_{\beth};\ \Gamma_\text{and};\
\Phi_\text{sub};\ H_0;\ S_{n:n};\ \Omega_0 \rangle,\]


is the subcritical, zero-winding, zero-chirality resting state that transcendental
structure is driven away from. This is not a surprising answer --- it is the
\emph{expected} answer, which is exactly what a passing check looks like. $\checkmark$

\textbf{Directed distance.} Moving up from the polynomial algebra to EML requires
acquiring $\Phi_c$, $\Omega_Z$, and $H_1$ --- three independent structural
conditions, each costly. Moving down only requires dropping them.
$d_{\to}(\text{polynomial}, \text{eml}) \gg d_{\to}(\text{eml}, \text{polynomial})$:
the direction of strain points the right way. $\checkmark$

\textbf{Multi-session convergence.} Three coordinates were derived by independent
arguments within this document and converged: $\Phi_c$ from the snapping manifold
definition in \pilcrow{}I and from the basin analysis in \pilcrow{}VI; $\Omega_Z$ from the branch
structure in \pilcrow{}I and from $\pi_1(\mathbb{C}^\times) = \mathbb{Z}$ in \pilcrow{}VII; $P_{\pm}$
from the adjoint structure of subtraction in \pilcrow{}II and from the external
$\mathbb{Z}_2$ orbit in \pilcrow{}IV. Then, immediately after completing \pilcrow{}X, a second
derivation was performed from scratch --- each primitive assigned independently,
without reference to \pilcrow{}I. The result:


\[\langle D_{\infty};\ T_{\bowtie};\ R_\text{dagger};\ P_{\pm};\ F_{\hbar};\
K_\text{slow};\ G_{\aleph};\ \Gamma_\text{seq};\ \Phi_c;\ H_1;\ S_{1:1};\
\Omega_Z \rangle.\]


Distance from session one: $d = 0$. This is the check that was flagged as open
in the first draft of this section and has since been closed. $\checkmark$

\textbf{Self-encoding test.} The EML compiler applied to $\text{eml}(x,y)$ returns
$\text{eml}(x,y)$ at depth 1. The EML representation of EML is EML. There is
something worth pausing on here: a system that asserts the exp/ln duality without
proving it can nonetheless represent its own structure exactly, because
self-representation does not require the Frobenius condition --- it only requires
that the grammar reach depth 1, which $\text{eml}(x,y)$ trivially satisfies.
This is the precise difference between $O_2^\dagger$ (can represent itself) and
$O_{\infty}$ (encoding and decoding are mutually inverse). The self-encoding test
passes; the self-inversion test does not. Both facts are correct. $\checkmark$

\textbf{Prediction consistency.} Four predictions from the encoding are falsifiable
against known results: (a) no EML tree generates a Stark unit --- consistent with
algebraic number theory, since Stark units require $P_{\pm}^{\text{sym}}$ and the
$P$ bottleneck prevents EML from reaching it; (b) snapping success rate decays
exponentially with depth --- confirmed by the paper''s data (100\%, $\sim 25\%$,
$<1\%$, $0\%$ at depths 2, 3--4, 5, 6); (c) no real-only EML-type Sheffer exists
--- confirmed by the paper''s negative search; (d) the ternary $T$ requires a
distinguished constant --- confirmed by $T(x,x,z) = \log_x z$. No predictions are
falsified. $\checkmark$

\textbf{The obligation discharged.} All nine checks pass. ENCODING\_EPISTEMOLOGY
states that an encoding satisfying all checks simultaneously, with no independent
derivation producing a persistent counter-encoding, is ''''not a choice --- it is a
finding.'''' The second session produced $d = 0$. The type assignment is a finding.

What it finds is this: EML is the highest sub-Frobenius type in the elementary
function algebra, and the gap between $P_{\pm}$ and $P_{\pm}^{\text{sym}}$ is not
a deficiency in the operator but a structural fact about what the elementary
functions can and cannot internally prove about themselves.

\begin{center}
\rule{0.5\textwidth}{0.4pt}
\end{center}

\emph{Structural encoding: $\langle D_{\infty};\ T_{\bowtie};\ R_\text{dagger};\ P_{\pm};\
F_{\hbar};\ K_\text{slow};\ G_{\aleph};\ \Gamma_\text{seq};\ \Phi_c;\ H_1;\
S_{1:1};\ \Omega_Z \rangle$. Tier: $O_2^\dagger$.}


\end{document}
